% !PRINTED 80
\weberpage{80}
%<<CONT

apart from sign, agrees with the subdeterminant of the \((n-\nu)\)-th order
\[
        A_{\alpha_{\nu+1}\ldots\alpha_n}^{\beta_{\nu+1}\ldots\beta_n}.
\]
Thus the complex of terms (6) is denoted by
\begin{equation}
        A_{\alpha_1,\alpha_2,\ldots,\alpha_\nu}^{\beta_1,\beta_2,\ldots,\beta_\nu}
        B_{\alpha_1,\alpha_2,\ldots,\alpha_\nu}^{\beta_1,\beta_2,\ldots,\beta_\nu}.
\tag{7}
\end{equation}

If now, for \(\alpha_1,\alpha_2,\ldots,\alpha_\nu\), one chooses each combination of \(\nu\) of the digits \(1,2,\ldots,n\), whose number (according to \S~7) is \(B_\nu^{(n)}\), while \(\beta_1,\beta_2,\ldots,\beta_\nu\) is held fixed, one obtains the same number of complexes of the form (7), and every term of the determinant \(A\) occurs in one and only one of them.

\gesperrt{V. Consequently, by summing all expressions (7), one obtains the determinant \(A\):}
\begin{equation}
        A=\sum^{\alpha}
        A_{\alpha_1,\alpha_2,\ldots,\alpha_\nu}^{\beta_1,\beta_2,\ldots,\beta_\nu}
        B_{\alpha_1,\alpha_2,\ldots,\alpha_\nu}^{\beta_1,\beta_2,\ldots,\beta_\nu}.
\tag{8}
\end{equation}
Of course one may also keep the combination of the \(\alpha\)'s fixed and sum with respect to the \(\beta\)'s.

This is Laplace's theorem.

Another representation of the determinant \(A\) by the first and second subdeterminants is obtained as follows.

Choose in \(A\) a row and a column which intersect in an element, say \(a_\nu^{(\mu)}\). In every term of \(A\) there occurs one element with lower index \(\nu\) and one element with upper index \(\mu\). Thus in \(A\) we have first the complex \(a_\nu^{(\mu)}A_\nu^{(\mu)}\), and furthermore the various complexes \(a_\nu^{(i)}a_k^{(\mu)}A_{\nu,k}^{i,\mu}\), in which \(i\) may denote any index different from \(\mu\), and \(k\) any index different from \(\nu\).

\gesperrt{VI. We can therefore put:}
\begin{equation}
        A=a_\nu^{(\mu)}A_\nu^{(\mu)}+
        \sum^{i,k} a_\nu^{(i)}a_k^{(\mu)}A_{\nu,k}^{i,\mu},
\tag{9}
\end{equation}
or, according to IV,
\begin{equation}
        A=a_\nu^{(\mu)}A_\nu^{(\mu)}-
        \sum^{i,k} a_\nu^{(i)}a_k^{(\mu)}A_{\nu,k}^{\mu,i}.
\tag{10}
\end{equation}
We further observe concerning this theorem that \(A_{\nu,k}^{\mu,i}\) is the \gesperrt{first subdeterminant}, corresponding to the element \(a_k^{(i)}\), of the \((n-1)\)-rowed determinant \(A_\nu^{(\mu)}\); for \(A_{\nu,k}^{\mu,i}\) is the coefficient
%CONT>>
